\section{The coarse error on the selected family member}

% BEGIN J79H1
\paragraph{J79H1. The coarse invariant bounds at constant 120.}
Fix the family member, its good points and its hub $X$ from F78D2.
Let $i$ satisfy all three guards in the coarse invariant condition:
it is twisted, its residue prime $\ell$ lies outside the two type
primes, and it is in the indicated range relative to $X.\mathrm{exc}$.
For the first two bounds, C78B3 and C78B4 give
$c_B(i)\le120tn n_V H$ and $h_0(i,C_i)\le120t$,
where $H=\mathrm{large}_2(L)$.
The third bound is an implication at a split prime. In that case,
for each declared $u\in C_i$, C78B3 gives
$c_0(i,u)\le120n n_VH$. Fix any good point $\pi$ at $i$. S79P1
identifies its twisted lattice $T$ over the coefficient field of
the source bridge. S79P2 proves that the associated representation
$V$ is crystalline. S79P3 and S79P4 give the same bound $h$ for
both ends of its jump interval and exclude zero at every embedding.

Apply W79H5 with $M_0=h_0=h$ on this layer of the chosen functor.
In A79L6--A79L8 use the de Rham comparison of W79H1 for $V$.
The good coefficient field contains every embedding of the local
field, so choose one of them, $\tau$. Write $a_\tau$ for the
exponent of A79L5 computed with the distinct jumps at $\tau$.
For every invariant class $v$ in $V/T$, A79L8 gives
$\ell^{a_\tau}v=0$. S79P5 gives $a_\tau\le c_0(i,u)$, whence
\[
       \ell^{c_0(i,u)}v=\ell^{c_0(i,u)-a_\tau}\ell^{a_\tau}v=0 .
\]
The injective map in C78B5 carries the original invariant classes
to these local invariants. Its normalized-length calculation now gives
$h_W(i,u,\pi)\le r c_0(i,u)\le120rn n_VH$.
This proves the split implication for every column and good point;
it is vacuous if there is no column or no good point.
These inequalities are precisely the three clauses of
$\mathrm{InvariantsBoundsWithCoarseColumnConstant}(L,X,120)$.
% END J79H1

% BEGIN J79H2
\paragraph{J79H2. Ten conclusion fields on one hub at one cutoff.}
Assume the original $\mathrm{Hypotheses}_2$, the family assumptions
$\mathrm{TwGood}$ and $\mathrm{FactsF}$ of F78D2, and $C_2\ge1$.
Take the single constant $C_s$ and cutoff $x_0$ of J78H1. For each
$x\ge x_0$ above the family threshold, and each allowed family
index there satisfying $\mathrm{ShaFinOdd}$, take the hub chosen
in J78H1 using the arithmetic exceptional set of that member.
It satisfies simultaneously
\[
 \begin{gathered}
 \mathrm{pos},\quad\mathrm{pass},\quad\mathrm{mult},\quad
 \mathrm{rank},\quad\mathrm{kill},\quad\mathrm{num},\\
 \mathrm{bb},\quad\mathrm{exc},\quad\mathrm{par},\quad\mathrm{small}.
 \end{gathered}
\]
Indeed its canonical data give the split and inert error identities
of C78B2. Apply that result to J79H1 with $c=120$ to obtain
$\mathrm{BBLe}(L,X,7200)$, since $60\cdot120=7200$.
The coefficient of this constant in the bound is
$fnt^2 n_V H^2\ge0$. As $C_s\ge22365\ge7200$, the same
inequality gives $\mathrm{BBLe}(L,X,C_s)$, the field
$\mathrm{bb}$. Everything else in the display holds on this $X$
by J78H1. In particular the assertion retains the same good-point
sets, decomposition, exceptional primes and family cutoff in all ten
fields. The other fields of the original conclusion are
$\mathrm{mid},\mathrm{large},\mathrm{two},\mathrm{ram},\mathrm{typ}$.
% END J79H2
